\section*{Conclusion}
\addcontentsline{toc}{section}{Conclusion}

Le mémoire a reconstruit sur des données françaises le modèle G2++ de \citet{BrigoMercurio2006}, dans la lecture
de \citet{Russo} où le second facteur est le spread, pour évaluer une option vanille collatéralisée sur l'OAT
4,10~\% 2046 dont le forward est fixé par le repo. Le modèle est recentré sur ce forward comme il est recalé sur
les courbes, et l'option a un prix exact, validé par un Monte-Carlo sans discrétisation.

\subsection*{Réponses aux hypothèses}

L'hypothèse H1 est rejetée. La courbe OAT étant recalée par construction, la procédure de calibration de
\citet{Russo} ne mesure que la convexité de la courbe ; son optimum exact ($\rho = -1$, $a_y = a_x$,
$\sigma_y = \sigma_x$) rend le taux OAT déterministe, et elle ne peut identifier aucun paramètre du spread.

L'hypothèse H2 est confirmée pour la mesure de pricing et rejetée pour la mesure historique. La structure par
terme des volatilités du spread, comme la coupe des spreads dans le filtre de Kalman, identifie une vitesse de
retour risque-neutre quasi nulle ; la GMM estime $\sigma_y$ à 33,7~pb par an et $\rho$ à 0,13, le filtre de Kalman
33,1~pb et 0,15. Mais $\rho$ varie de 0,06 à 0,43 selon le tenor et la fenêtre, bien au-delà de son écart-type, six
ans de données ne suffisent pas à identifier la dynamique historique, et les tests de Hansen et de Ljung-Box
rejettent le spread à un facteur. Les deux régimes de volatilité, de 23 et 51~pb par an, sont trop courts pour
peser sur le prix.

L'hypothèse H3 est confirmée : le facteur spread renchérit le call à la monnaie de 17~\% à toutes les échéances,
et les volatilités réalisées des taux à vingt ans confirment cet ordre de grandeur (rapport OAT~/~€STR de 1,17 dans
les données et de 1,18 dans le modèle).

L'hypothèse H4 est rejetée avec les paramètres estimés. Les deux vitesses de retour étant quasi nulles, un
Hull-White à un facteur sur la courbe OAT, à la volatilité totale du taux OAT, reproduit le call à la monnaie à
1,9~\% près, et, à volatilité ajustée, les prix de 90 à 110~\% du forward à 0,05~\% près. Ce qui compte pour une
option sur une OAT est le niveau de cette volatilité, supérieur de 15~\% à celui que donnent les swaptions €STR ;
sa mesure demande l'historique du spread, faute d'options sur OAT.

L'hypothèse H5 est confirmée à la monnaie et rejetée en dehors. L'approximation à poids gelés ne s'écarte que de
0,12~\% du prix exact à la monnaie lorsque les poids sont gelés à leur valeur forward recentrée ; en dehors, elle
coûte jusqu'à $-3{,}0\,\%$ sur le call à 110~\% et $+4{,}3\,\%$ sur le put à 90~\%, à un an. Le pricer utilise donc le
prix exact.

À la question posée en introduction, la réponse est donc nuancée. Il faut mesurer la volatilité du spread, qui
renchérit l'option de 17~\%, et les séries historiques permettent de l'identifier sous la mesure de pricing. Mais
pour une option sur une seule OAT, un facteur de volatilité augmentée suffit à la représenter. Enfin, aux longues
échéances, le prix dépend autant du niveau du financement en repo que du facteur spread : à dix ans, translater la
courbe repo de $-25$ à $+25$~pb déplace le call à la monnaie de $-21\,\%$ à $+24\,\%$.

\subsection*{Limites}

La première limite est la courbe repo, fictive faute de cotations. Les options sont limitées à dix ans pour que
cette courbe couvre toutes les échéances sans extrapolation ; une courbe de marché à ces durées reposerait
elle-même sur un repo à terme peu liquide. Le recentrage $\delta$ place par ailleurs dans la moyenne du prix une
perte anticipée que sa dispersion ignore. L'approximation reste limitée tant que $\delta$ est petit, 33~pb au plus
ici, mais elle grandirait avec l'écart entre le spread repo et le spread de l'OAT.

Le modèle ne comporte qu'un facteur par courbe, et les tests de Hansen et de Ljung-Box rejettent cette structure :
un seul facteur ne reproduit ni les déformations de pente du spread (deuxième composante principale, 8~\% de la
variance), ni des corrélations qui varient d'un tenor à l'autre. Le spread gaussien peut devenir négatif, ce qui
n'est pas absurde pour un spread OAT–€STR, négatif à deux ans presque tout au long de 2021 et 2022. Seule une
composante de crédit devrait rester positive ; un processus affine positif \citep[cours~1]{MPR2025} l'imposerait.
Le spread, qui mêle crédit, liquidité et effets d'offre, est traité comme un facteur de risque de marché : le modèle
ne décrit pas un événement de crédit sur l'État français.

La calibration est mixte : la volatilité des taux est implicite, celles du spread et la corrélation sont
historiques. Les deux sources sont compatibles sur 2021-2026 pour les taux, sans garantie qu'elles le restent, et
pour le spread l'égalité des volatilités sous les deux mesures reste une hypothèse. La volatilité des taux est
constante ; pour une option européenne, seule compte la variance intégrée jusqu'à l'échéance, et l'effet d'une
volatilité dépendant du temps est borné par les résidus de la diagonale, 2,3~pb au plus. Les conventions sont
simplifiées (base ACT/365 sans calendrier, tenors intermédiaires interpolés, rendements génériques de pair pris
comme taux zéro). Enfin, le modèle gaussien impose sa forme aux prix hors de la monnaie, qu'aucun paramètre ne
permet d'ajuster \citep[section~6.3.5]{Chaix2021}.

\subsection*{Prolongements}

Le premier prolongement est de remplacer la courbe repo fictive par une courbe de marché, au même format. Les
résidus du filtre de Kalman réclament un facteur de pente par courbe et des erreurs de mesure autocorrélées.
L'instabilité de $\rho$, que les régimes de volatilité n'expliquent pas, appellerait une corrélation stochastique,
par exemple un processus autorégressif de Wishart \citep[cours~4, slides~115-118]{MPR2025}, et un G2++ à
changement de régime, qui reste affine (cours~4, slides~44-45), traiterait des options plus courtes que les
régimes. On pourrait aussi séparer une composante de crédit positive et une composante de liquidité, comme
\citet{MonfortRenne2014}, ou introduire une prime de risque de variance sur le spread (cours~6, slides~116-118),
identifiable si des options sur OAT devenaient liquides. Enfin, la structure à deux facteurs compterait vraiment
pour des produits sensibles à l'écart entre les deux courbes, comme les options sur spread ou sur plusieurs titres.

\subsection*{Recul sur la démarche}

Trois enseignements dépassent le cas étudié. Le premier est de tester l'identification avant d'estimer : la
procédure de calibration du papier de référence paraissait naturelle, mais une ligne de calcul montre qu'elle ne
mesure que la convexité de la courbe, et sa mise en œuvre aurait produit des paramètres sans contenu. Le deuxième
est que la définition du contrat compte autant que le modèle : le forward repo, imposé par la pratique de marché,
déplace le prix à dix ans autant que le second facteur. Le troisième est qu'un modèle plus riche n'est utile que
si le produit est sensible à ce qu'il ajoute ; ici, le second facteur se résume, pour le prix, à un supplément de
volatilité, et c'est la mesure de ce supplément qui demande le plus de travail économétrique.
\acompleter{apport du stage pour l'élève et prise de recul personnelle, en quelques phrases}.
